\section{Related Work}
\label{sec:related}

The concept of a design fingerprint --- a characteristic multi-dimensional signature that
identifies an artifact's distinctive pattern of feature emphasis --- has been applied
in several domains. In music, timbre fingerprints derived from MFCC features identify
instrument signatures independent of pitch and loudness~\cite{aucouturier2003music}.
In visual art, style fingerprints derived from CNN feature maps identify artist signatures
independent of subject matter~\cite{gatys2016image}. The common structure is within-item
normalization: subtracting a global quality signal to expose relative emphasis.

In game design, the closest prior work is BGG's ``family'' and ``mechanic'' tag system,
which creates a form of discrete fingerprint from community-assigned category labels.
Our approach differs in using continuous scored axes rather than binary tags, and in
computing fingerprints as the relative deviation from a game's own average rather than
as an absolute vector --- a distinction that becomes critical when comparing games across
quality tiers. The within-game z-score normalization follows the psychometric tradition
of ipsative scoring~\cite{meade2004psychometric} adapted for design analysis.
